
\chapter{Part V Review: Correlators, Diagrams, and Scattering}

Perturbation theory is a controlled expansion around the Gaussian integral.  Wick's theorem is the combinatorics of that Gaussian, and Feynman diagrams are a compact record of the resulting terms.

\section{A guided reconstruction}

% QFTFIGURE-BEGIN partreview05-01-part-map
\qftfigure{figures/instructional/partreview05/partreview05-01-part-map.pdf}{The map collects the part's main constructions into a visual route so the reader can see how the chapters support one another.}{fig:partreview05-01-part-map}
% QFTFIGURE-END partreview05-01-part-map












For \(\lambda\phi^4\), one vertex contributes \(-\ii\lambda\).  Propagators connect the vertex to other points or to other vertices.  LSZ then explains why poles of correlation functions are related to scattering amplitudes.  The slogan is short, but the derivation is a chain: source derivatives, contractions, Fourier transforms, pole residues, and phase-space normalization.

\begin{exercise}
List the steps that turn a four-point correlation function into a scattering amplitude.
\begin{answer}
Compute the connected correlator, Fourier transform, identify external poles, amputate the external propagators, put external momenta on shell, and include the correct normalization conventions.
\end{answer}
\end{exercise}
